\chapter{Introduction}
\label{chap:introduction}

The South African mining industry continues to experience a disproportionate burden of occupational lung disease, driven by chronic exposure to silica dust \citep{silicosis_elimination_sa_2020}, diesel particulates \citep{badenhorst2025diesel, diesel_pm_health_effects_2020}, and microbial bio-aerosols in poorly ventilated mine environments.  
Miners remain at particularly high risk for silicosis and tuberculosis (TB) \citep{silicosis_elimination_sa_2020}, as well as their co-occurrence, known as silico–TB.  
Radiographic screening is the primary triage tool for these diseases, yet many clinics operate with ageing X-ray systems that produce heterogeneous image quality and inconsistent contrasts \citep{toua2022chest, ehrlich2024silicotuberculosis}.  
Consequently, the available datasets \citep{chambon2024chexpert,nguyen2022vindr, cattamanchi2024coda} are fragmented across machines, hospitals, and annotation protocols—leaving deep learning models under-trained precisely on the complex edge cases that matter most: comorbid and atypical presentations.  
Developing data-efficient methods that can faithfully expand such underrepresented classes is thus essential for building robust and equitable diagnostic models.

\citet{predictors2020silicosis} demonstrate that despite long-standing awareness of occupational silica exposure risks, silicosis remains a persistent and widespread disease among South African gold miners. 
Their epidemiological findings reveal that cumulative exposure over time remains the dominant risk factor, and that many ex-miners continue to develop silicosis years after leaving the mining environment. 
Importantly, the study highlights the strong and recurrent co-morbidity between silicosis and tuberculosis (TB), particularly in aging mining populations, emphasizing the need for continued surveillance and improved medical data integration to understand disease progression and co-occurrence.

\citet{silicosis_elimination_sa_2020} similarly report that the goal of silicosis elimination in South Africa remains distant, hindered by inadequate dust control, unreliable quartz measurement, and limited long-term health monitoring. 
They caution that even exposure levels within current regulatory limits are insufficient to prevent disease onset and that silicosis substantially increases susceptibility to TB infection. 
The authors argue that effective disease mitigation will require sustained epidemiological follow-up and new data-driven approaches capable of characterizing how silica-related lung damage interacts with infectious pathology.

Together, these studies underscore the entrenched relationship between occupational silica exposure, silicosis, and TB co-morbidity.
While policy and public health interventions have primarily targeted exposure reduction, there remains a critical gap in quantitative and computational modeling of co-occurrence patterns and imaging phenotypes. 
Given that chest radiography remains the primary diagnostic tool in low-resource mining regions, yet annotated datasets are small and heterogeneous, developing generative models capable of synthesizing plausible Silico-TB images could support both data augmentation and scientific exploration of disease interactions.

% Addressing the diagnostic and data challenges posed by Silico-TB requires methods that can reason about and reconstruct the latent structure of medical images rather than rely solely on existing, incomplete datasets. 
% In recent years, deep learning has demonstrated strong potential for advancing medical imaging through pattern recognition, segmentation, and automated diagnosis. 
% However, the majority of AI applications remain fundamentally discriminative: they learn to classify or segment from fixed, labelled datasets without understanding the broader distributional structure of disease patterns. 
% Such approaches struggle in domains—like occupational lung disease—where data are sparse, heterogeneous, and under-represented for rare comorbidities. 
% To move beyond these limitations, there is growing interest in \emph{generative} methods that model the full data distribution itself, enabling the synthesis and reconstruction of images that remain statistically and anatomically coherent.

% Generative modeling has become a central paradigm in contemporary machine learning, offering a probabilistic framework for learning complex data distributions and synthesizing realistic samples from them.  
At its core, a generative model seeks to approximate the data distribution \(p_{\text{data}}(x)\) by learning a parameterized distribution \(p_\theta(x)\) from which new samples can be drawn that resemble those in the training data.  
Unlike discriminative approaches that predict conditional mappings \(p(y \mid x)\), generative models aim to reconstruct the underlying structure of data itself, enabling simulation, reconstruction, and inference under uncertainty \citep{bond2021deep,tomczak2024deep}.  
Recent advances \citep{rombach2022high, pmlr-v139-nichol21a, 2022arXiv220406125R} —most notably diffusion-based models—have shown remarkable success across domains ranging from natural image generation \citep{rombach2022high} to protein folding \citep{wu2024protein, skreta2024superposition}, establishing themselves as state-of-the-art methods for modeling high-dimensional distributions.

Diffusion models, first introduced by \citet{pmlr-v37-sohl-dickstein15} and
popularised in their denoising formulation by \citet{NEURIPS2020_4c5bcfec},
define a generative process by gradually corrupting data with Gaussian noise and
learning a reverse denoising process that reconstructs structure from noise,
either via discrete Markov chains (DDPMs) or continuous-time SDEs. Their
likelihood-based training, stability, and controllable sampling dynamics have
made them highly effective for scientific and medical imaging applications
\citep{accelerated_mri_motion_correction_2023,score_based_diffusion_mri_2022}.
Traditional conditional generative models require retraining when new attributes
or conditions arise, whereas \emph{composable diffusion models} enable
inference-time combination of pretrained components by viewing diffusion as an
energy-based system whose score approximates $\nabla_x \log p(x)$. This makes it
possible to perform logical operations such as conjunction and disjunction
\citep{NEURIPS2020_49856ed4,liu2022compositional} without retraining. Recent
methods such as \textsc{SuperDiff} \citep{skreta2024superposition} extend this
capability using an Itô-based density estimator that allows algebraic
combination of pretrained diffusion processes in a mathematically consistent
manner. Applying these compositional operators to disease-specific diffusion
models—such as those trained separately on silicosis and tuberculosis—enables
the synthesis of plausible co-morbid pathologies without dataset merging, making
composition an attractive framework for flexible, data-efficient medical image
generation.


Beyond data augmentation, controllable generative models open new avenues for scientific reasoning.  
Clinicians and researchers can perturb learned pathology factors—such as varying silica exposure severity while fixing TB lesion distribution—to test hypotheses about disease manifestation and progression.  
Such “in silico” experimentation provides an interpretable and reproducible medium for exploring diagnostic cues, aiding radiologists in understanding how visual features correlate with disease dynamics.  
These capabilities align with South Africa’s public health objectives by supporting remote triage, improving tele-radiology workflows in mining towns, and facilitating targeted screening strategies under data scarcity \citep{SAJR2647}.


This work aims to establish a principled framework for compositional diffusion
modelling, with an emphasis on understanding and generating clinically plausible
co-morbid chest pathologies. The central objective is to determine \emph{when}
and \emph{why} compositions of pretrained diffusion models produce meaningful
hybrid representations rather than incoherent or anatomically implausible
samples. The specific goals are:

\begin{itemize}
    \item \textbf{To characterise the structural requirements for successful
    compositionality.}  
    We investigate how latent alignment, local disentanglement of semantic
    factors, conditional factorisation, and log-density geometry govern the
    behaviour of score-based, energy-based, and density-based operators, and
    when these operators can be expected to preserve semantic intent.

    \item \textbf{To develop and evaluate guidance mechanisms that stabilise
    compositional sampling.}  
    We examine whether continuous properties—such as severity, localisation,
    or morphological attributes—can be used to regularise and control hybrid
    generations, improving robustness and reducing hallucination under
    distribution shift.

    \item \textbf{To analyse geometric compatibility between disease-specific
    generative models.}  
    Using tools from Riemannian diffusion and latent manifold analysis, we
    assess whether the generative manifolds of conditions such as TB and
    silicosis support coherent composition, and whether geometric mismatch
    explains characteristic failure modes observed in practice.
\end{itemize}

% By addressing the underrepresentation of co-morbid chest pathologies within
% South Africa’s mining sector through a mathematically grounded and
% data-efficient generative framework, this research aims to bridge theoretical
% advances in diffusion modelling with the practical needs of occupational health
% imaging in resource-constrained environments. Co-morbid conditions such as
% Silico-TB constitute a scientifically meaningful testbed, as their non-linear
% and interdependent pathology patterns expose precisely the structural
% difficulties that current compositional methods are not equipped to handle.

In support of this investigation, the proposal is structured as follows.  
Chapter~\ref{chap:background} presents the core background, covering diffusion
models, their mathematical formulations, and the main families of compositional
operators, together with recent advances in context-guided and Riemannian
diffusion.  
Chapter~\ref{chap:related_works} surveys generative modelling in medical
imaging, including synthetic radiography and prior approaches to
compositionality.  
Chapter~\ref{chap:research_objectives} introduces the motivation, research
goals, and hypotheses that structure the investigation.  
Chapter~\ref{chap:research_method} then outlines the methodological
framework, detailing the three experimental phases spanning controlled medical
testbeds, property-guided composition, and geometric analysis using Riemannian
diffusion.  
Finally, Chapter~\ref{chap:conclusion} summarises the anticipated contributions
and the broader scientific and clinical significance of the work.


